\documentclass{article}
\usepackage{arxiv}

\usepackage[utf8]{inputenc}
\usepackage[T1]{fontenc}
\usepackage{hyperref}
\usepackage{url}
\usepackage{booktabs}
\usepackage{amsfonts}
\usepackage{amsmath}
\usepackage{nicefrac}
\usepackage{microtype}
\usepackage{graphicx}
\usepackage{natbib}
\usepackage{doi}
\usepackage{cleveref}
\usepackage{xcolor}
\usepackage{multirow}
\usepackage{array}
\usepackage{caption}
\usepackage{subcaption}
\usepackage{float}

\graphicspath{{figures/}}

% Allow floats to share pages with text more readily (reduces large gaps).
\renewcommand{\topfraction}{0.9}
\renewcommand{\bottomfraction}{0.8}
\renewcommand{\textfraction}{0.07}
\renewcommand{\floatpagefraction}{0.7}
\setcounter{topnumber}{3}
\setcounter{bottomnumber}{2}
\setcounter{totalnumber}{4}

\title{ReGainX: Adaptive Reinforcement Learning Exoskeletons\\
for Real-Time Mitigation of Parkinson's Motor Impairments}

\author{%
  Rohan Bankapur \\
  Independent Researcher \\
  \texttt{Rohannb02@gmail.com} \\
  Code: \href{https://github.com/TerminalRocketship45/ReGainX}%
  {\texttt{github.com/TerminalRocketship45/ReGainX}}
}

\begin{document}
\maketitle

%------------------------------------------------------------------
\begin{abstract}
Parkinson's disease (PD) limits daily elbow movement in over 10 million
people, causing progressive muscular degeneration and idiosyncratic
bradykinesia. Existing exoskeletons fail to generalise to Parkinson's
patients due to a lack of representative data, hardware limitations, and
calibration difficulty. ReGainX introduces a three-stage multi-agent
reinforcement learning framework that creates diverse Parkinson's digital
twins using Proximal Policy Optimisation (PPO). A custom RecurrentPPO
algorithm with an LSTM policy fuses EMG and IMU data for real-time,
adaptive torque control. The trained policy achieves a mean trajectory
match of $r{=}0.494$ and a goal achievement rate of $88.8\%$ across
4\,000 stratified evaluation episodes spanning the full PD severity
spectrum, without per-session recalibration. An ablation study demonstrates that the combined-trained policy achieves
a $35.7\%$ higher trajectory match than exoskeletons trained on muscle
weakness alone, sustaining a mean reward of $382.30$ at the most severe
impairment level versus $9.26$ for the degeneration-only
policy---validating combined impairment modelling as a necessary
condition for cross-patient generalisability.
\end{abstract}

\keywords{Parkinson's disease \and exoskeleton \and reinforcement learning
\and musculoskeletal simulation \and bradykinesia \and multi-agent systems}

%==================================================================
\section{Introduction}
\label{sec:intro}

Parkinson's disease (PD) affects over 10 million people globally, costing
\$82.2 billion annually in the U.S.~\citep{chaudhuri2025, parkinson2026}.
Two compounding motor impairments progressively restrict daily elbow
movement: \emph{muscular degeneration}, the progressive weakening of
muscle force by sarcopenia and fatigue, and \emph{bradykinesia}, the
fluctuating slowing and delay of voluntary movement. Critically, PD motor
severity does not remain fixed---patients cycle between ON and OFF states
within a single day as medication efficacy fluctuates, producing
unpredictable swings in impairment magnitude~\citep{ramesh2022}. Despite
this, even mildly affected patients show a $22\%$ force production deficit
at the elbow~\citep{ingram2021}. Existing wearable technology fails PD
patients in three ways: per-session manual calibration, unreliable noisy
sensing, and a lack of multimodal PD-specific datasets needed to train
generalisable devices~\citep{researchgate2023, arxiv2025}. Despite a
$90\%$ interest in medical wearables, only $8\%$ of PD patients adopt
them~\citep{hirczy2024}.

Elbow exoskeletons have evolved significantly since early devices relied
on rigid trajectory control, where the controller moved the joint through
a fixed path regardless of user effort---a system that becomes inflexible
when patient states constantly change~\citep{gasperina2021, myoelectric2022}.
As reviewed by \citet{supriyono2025}, the elbow joint is the most targeted
joint in upper-limb exoskeleton design, yet control strategies have
historically required patient-specific and session-specific parameter
tuning that cannot account for short- and long-term severity fluctuations.
The field has progressively shifted toward intent-detection using
myoelectric signals; EMG captures continuous, multi-channel neural drive
data reflecting both movement intention and muscle activation state, but
surface EMG remains highly susceptible to noise and session-to-session
variability---barriers compounded in PD patients whose neural drive
signals carry abnormal bursting patterns from basal ganglia
dysfunction~\citep{supriyono2025, mdpi2025}.

Deep reinforcement learning (DRL) addresses this directly: rather than
calibrating to a fixed patient profile, a DRL policy trained across a
diverse distribution of patient conditions learns to generalise without
manual re-tuning between sessions. \citet{luo2024} validated this
principle when training a hip exoskeleton entirely in simulation and
deploying it to real hardware without a single human experiment, achieving
a $24.3\%$ reduction in metabolic walking cost. \citet{endrei2025} applied
DRL specifically to PD, demonstrating tremor suppression at the elbow and
shoulder across dynamic movements. Despite this, no prior work has trained
an autonomous elbow exoskeleton controller for muscular degeneration and
bradykinesia---the defined characteristics of PD motor impairments. The
absence of real multimodal PD patient datasets makes simulation-based
training the only viable path, and no existing framework has constructed a
proper environment to do this for exoskeletons.

ReGainX addresses this with a three-phase multi-agent DRL framework.
Using MyoSuite's high-fidelity musculoskeletal simulator~\citep{caggiano2022},
we first train a healthy digital twin of the human elbow with DRL, then
simultaneously apply randomised severities of muscular degeneration and
bradykinesia to generate a distribution of PD digital twins, and finally
train a DRL exoskeleton agent against the frozen PD controller across
approximately 20\,000 distinct PD patient profiles. The objective of this
paper is to assess whether deep reinforcement learning can produce an
autonomous elbow exoskeleton controller that provides real-time, adaptive
mitigation of motor impairments at the elbow across diverse PD patients.

%==================================================================
\section{Related Work}
\label{sec:related}

\subsection{Control Strategies for Upper-Limb Exoskeletons}

Elbow exoskeletons have evolved across three distinct control generations.
First-generation devices applied forces independently of voluntary human
activity---an approach fundamentally unsuited to patients with non-constant
impairment levels~\citep{gasperina2021}. The logical successor, intent-detection
through surface EMG, allows devices to respond to neural drive signals.
\citet{supriyono2025} confirm that EMG and IMU sensing now dominate the
field precisely because of their low cost and high information density.
However, the same review identifies sEMG noise, electrode placement shifts,
and IMU drift as the primary barriers preventing reliable cross-patient
generalisation~\citep{supriyono2025}.

\subsection{Exoskeleton Research in Parkinson's Disease}

Despite the scale of PD as a motor disorder, published exoskeleton work
targeting PD patients remains strikingly sparse. \citet{fortunati2024}
conducted a systematic review of powered lower-limb exoskeletons in PD
gait training and found only five original studies involving 46 patients
total---all stable on medication, all at Hoehn and Yahr stages below 3.
None used adaptive or learning-based control. Furthermore, all five
studies addressed gait, with zero coverage of upper limb function. The
elbow---a major site of PD motor decline---has not been the target of
any published PD-specific exoskeleton study to date, and there is
currently no standardised assessment or targeted intervention for upper
limb function in PD rehabilitation~\citep{capato2023}.

\subsection{Deep Reinforcement Learning for Exoskeleton Control}

The recognised limitation of static, patient-specific controllers has
motivated a shift toward DRL. \citet{luo2023} presented the first DRL
controller for a lower-limb rehabilitation exoskeleton trained on a
coupled musculoskeletal model, using domain randomisation of muscle
strength across quadriplegic, muscle-weakness, and hemiplegic conditions.
\citet{luo2024} extended this to experiment-free deployment.
\citet{endrei2025} applied DRL specifically to PD tremor suppression.
However, none of these works model the simultaneous interaction of
muscular degeneration and bradykinesia, nor construct a multi-agent
environment where a simulated patient and an exoskeleton agent
co-exist during training.

\subsection{Musculoskeletal Simulation as a Training Foundation}

The absence of PD datasets makes high-fidelity simulation the most viable
training framework. \citet{caggiano2022} introduced MyoSuite, a
contact-rich musculoskeletal simulation suite built on MuJoCo that
supports the modelling of physiological impairments---including fatigue
and sarcopenia---as native non-stationarities. \citet{wang2022} extended
this with MyoSim, operating over two orders of magnitude faster than
prior tools, making large episode counts computationally feasible.
No prior work has applied these capabilities to train a PD-specific
elbow exoskeleton controller under combined muscular degeneration and
bradykinesia.

%==================================================================
\section{Parkinson's Disease Simulation}
\label{sec:pdsim}

\subsection{Musculoskeletal Model and Joint Scope}

ReGainX's digital twins are built within MyoSuite's \texttt{myoElbow}
environment, a 6-muscle elbow model embedded in MuJoCo's physics engine,
comprising three primary flexors---BIClong, BICshort, and BRA---and three
primary extensors---TRIlong, TRIlat, and TRImed~\citep{caggiano2022}.
Muscle force generation follows Hill-type muscle-tendon dynamics, in which
each muscle's contractile force depends on its activation level, fiber
operating length, and contraction velocity~\citep{zajac1989}. The
resulting joint torque is expressed through the rigid-body equation of
motion:
\begin{equation}
  \boldsymbol{\tau} = M(\mathbf{q})\ddot{\mathbf{q}}
    + C(\mathbf{q}, \dot{\mathbf{q}})\dot{\mathbf{q}}
    + G(\mathbf{q}),
\end{equation}
where $M(\mathbf{q})$ is the inertia matrix, $C(\mathbf{q},\dot{\mathbf{q}})$
captures Coriolis and centrifugal forces, and $G(\mathbf{q})$ is the
gravitational torque vector. The simulated joint operates over $0$ to
$2.4\,\mathrm{rad}$ ($\approx 138^\circ$), consistent with the clinically
established active range of motion~\citep{cools2017, morrey1981}. The
exoskeleton provides a normalised torque $\tau_\text{exo} \in [0,1]$,
representing a fraction of the biological joint torque capacity
following~\citet{durandau2019}.

The elbow is targeted for two clinical reasons: elbow extensors show
the largest force production deficit ($22\%$) even in mildly affected,
medicated PD patients~\citep{ingram2021}, and the hand-to-mouth
movement---essential for eating and drinking---is significantly altered in
PD through reduced elbow flexion velocity and range of motion,
correlating with UPDRS-III scores~\citep{caligiuri2018}.

\subsection{Muscular Degeneration Model}

Muscular degeneration in PD arises from sarcopenia (OR $3.98$ vs.\
controls, prevalence $10.9$--$31.4\%$) and chronic muscle fatigue
(pooled prevalence $\approx 50\%$)~\citep{hart2023, ponsoni2023, siciliano2024}.
ReGainX adopts the combined degeneration equation from the MyoSuite
framework~\citep{caggiano2022}:
\begin{equation}
  F_\text{deg} = \alpha_1 \cdot F_\text{max}
    \cdot e^{-(\alpha_2 + \beta t)}
    = m_\text{loss} \cdot m_\text{fati},
\end{equation}
where $\alpha_1 \in [0,1]$ is the sarcopenic mass-loss factor,
$\alpha_2 \in [0,1]$ is the fatigue onset parameter (weighted toward
higher values reflecting patients who require assistance), and $\beta$
is held at MyoSuite's default value from the Cheema et al.\ (2020)
three-compartment fatigue model. All parameters are independently
randomised per episode.

\subsection{Bradykinesia Model}

Bradykinesia produces two quantifiable deficits: a delay in movement
initiation and a reduction in movement speed and amplitude. Off-medication
PD patients move $41\%$ slower than controls with $29\%$ shorter
displacement~\citep{agostino2015}. ReGainX models bradykinesia as:
\begin{equation}
  N_\text{brady} = x(i - \Delta i) \cdot \alpha,
\end{equation}
where $\Delta i \in [1.1, 1.4]$ is a temporal scaling factor (an action
taking $1.0\,\mathrm{s}$ healthy takes $1.1$--$1.4\,\mathrm{s}$ in the
PD patient) and $\alpha \in [0.3, 1.0]$ scales the movement amplitude,
reproducing the $29$--$41\%$ reduction documented
clinically~\citep{agostino2015, raethjen2018}.

\subsection{Combined PD Severity and Episode Diversity}

Both impairments are applied simultaneously across all training episodes:
the patient model has degraded peak force capacity \emph{and} delayed,
attenuated movement. The composite PD severity of each episode is
parameterised across both impairment scales, with start angle and target
goal angle randomised uniformly over $[0, 2.4]$ rad. The policy
encounters no two identical patient profiles during training and must
generalise across the full clinical spectrum.

Figure~\ref{fig:pd_modeling} illustrates the effect of each impairment
on muscle activation and joint trajectory.

\begin{figure}[tbp]
  \centering
  \includegraphics[width=\linewidth]{fig2_pd_modeling.pdf}
  \caption{Parkinsonian motor-deficit model. \textbf{Left:} muscle
    activation is reduced by degeneration and further delayed under
    combined PD impairment. \textbf{Right:} joint-angle trajectories
    show reduced reach and delayed onset; the vertical dashed line marks
    the akinetic delay $\Delta i$.}
  \label{fig:pd_modeling}
\end{figure}

%==================================================================
\section{Methodology}
\label{sec:method}

\subsection{Safety Design}
\label{sec:safety}

The MyoSuite exoskeleton model incorporates an upper-arm attachment
($0.101\,\mathrm{kg}$) and a forearm attachment ($0.111\,\mathrm{kg}$),
for a combined simulated weight of $0.212\,\mathrm{kg}$---within the
$0.2$--$0.5\,\mathrm{kg}$ range identified as feasible for patients with
reduced arm strength~\citep{pmc11852623, pmc12451468}. A healthy elbow
controller evaluated with the exoskeleton mass attached achieved perfect
trajectory match, confirming the modelled weight does not interfere with
voluntary movement.

Two explicit safety constraints are enforced by the environment. First,
the joint angle is hard-capped at $2.4\,\mathrm{rad}$, the anatomically
established upper limit of active elbow flexion~\citep{cools2017}.
Second, exoskeleton torque is normalised to $[0,1]$ as a fraction of
biological joint torque; the exoskeleton cannot deliver more torque than
the joint's own biological capacity.

\subsection{Multi-Agent Training Framework}
\label{sec:marl}

ReGainX implements a three-phase pipeline (Figure~\ref{fig:pipeline}),
sequentially developing a healthy musculoskeletal controller, a simulated
PD patient, and an exoskeleton agent. The patient controller is frozen
before exoskeleton training begins to prevent co-adaptation: if both
agents trained simultaneously, the patient model could learn to move in
unnatural patterns that inflate the reward---a failure that would not
transfer to real patients whose movement patterns are not shaped by the
device. ReGainX is designed so the exoskeleton adapts to the patient,
not vice versa.

\begin{figure}[tbp]
  \centering
  \includegraphics[width=\linewidth]{fig1_pipeline.pdf}
  \caption{Three-phase ReGainX training pipeline, shown with MuJoCo/MyoSuite
    renders at a fixed $2.0\,\mathrm{rad}$ target. \textbf{Phase 1:} a PPO
    healthy controller (trained) drives the six elbow muscles and reaches the
    target. \textbf{Phase 2:} the same controller is frozen and run under
    combined bradykinesia and muscular degeneration, stalling well short of
    the target. \textbf{Phase 3:} a RecurrentPPO exoskeleton agent (trained)
    supplies assistive torque to the frozen patient, restoring movement to the
    target. The exoskeleton is trained under domain randomisation over
    $\approx$20\,000 PD profiles.}
  \label{fig:pipeline}
\end{figure}

\paragraph{Phase 1 --- Healthy Controller Training.}
A musculoskeletal elbow controller is trained on
\texttt{myoElbowPose1D6MRandom-v0}---a healthy 6-muscle elbow model---via
PPO~\citep{sb3} for 250\,000 timesteps, with policy updates every
2\,048 steps. The action space is the 6-dimensional neural drive signal.
The reward is given each timestep, scaled to how much closer the joint
moved toward the target angle. Both starting and target angles are drawn
uniformly from $[0, 2.4]\,\mathrm{rad}$ per episode.

\paragraph{Phase 2 --- Parkinson's Patient Controller Setup.}
The healthy controller is transferred unchanged to the PD-impaired
environment (Section~\ref{sec:pdsim}). Joint position errors increase,
targets are missed, and reduced peak force combined with slowed activation
dynamics creates a realistic impaired profile. Randomising impairment
parameters alongside starting and target angles per episode generates
$\approx$20\,000 distinct patient profiles---serving as a
simulation-based alternative to real PD datasets, which are
scarce~\citep{researchgate2023, arxiv2025}. The controller is then frozen,
representing a PD patient still attempting voluntary movement but unable
to complete it.

\paragraph{Phase 3 --- Exoskeleton Agent Training.}
With the PD controller frozen, the exoskeleton agent is initialised and
trained for 2\,000\,000 timesteps using RecurrentPPO. Both agents act
within the same environment each timestep. The exoskeleton agent receives
an 8-dimensional observation and outputs a single normalised torque
$\tau_\text{exo} \in [0,1]$.

The observation space satisfies the Markov property: it consists of the
6 current muscle activations from the patient controller, the current
joint angle $q$, and the current joint velocity $\dot{q}$. Velocity
encodes the rate of joint state change without requiring access to prior
states. The target angle and starting angle are excluded: it is not
feasible to input precise joint angles every millisecond in deployment,
and it is impossible to extract the patient's intended target from neural
signals. Patient diversity is achieved through continuous domain
randomisation of \texttt{force\_scale}, \texttt{activation\_slowdown},
and muscle fatigue fraction at each episode.

\subsection{Policy Architecture}
\label{sec:policy}

The exoskeleton policy is a RecurrentPPO~\citep{sb3contrib}, which
augments the standard PPO actor-critic framework with an LSTM layer
between the observation input and the policy head. At each timestep,
the LSTM maintains a cell state and hidden state carried forward across
timesteps within an episode. At episode boundaries both states are
reset, ensuring the policy re-initialises its temporal context for each
new patient profile.

This formulates the control problem as a \emph{Partially Observable
Markov Decision Process} (POMDP): the observation deliberately excludes
the user's goal angle. The LSTM builds up evidence about the patient's
activation pattern over a movement, allowing the policy to infer the
direction of intended movement and the degree to which impairment will
prevent its completion. A feedforward policy acting only on the current
observation cannot distinguish a progressively fatiguing patient from one
at a stable impairment level.

A shorter rollout of $n_\text{steps}{=}256$ was used (vs.\ 2\,048 in
Phase 1) to produce more frequent LSTM state resets during training,
amplifying the episode boundary signal and bounding truncated
backpropagation through time to avoid vanishing gradients.

Training reward curves for all policies are shown in
Figure~\ref{fig:training_curves}.

\begin{figure}[tbp]
  \centering
  \includegraphics[width=\linewidth]{fig3_training_curves.pdf}
  \caption{Training reward curves (rolling mean, window~$=20$, with
    $\pm1\sigma$ band). The RecurrentPPO brady+deg policy converges
    substantially more slowly than MLP policies, consistent with the
    higher sample complexity of learning a useful LSTM hidden state.}
  \label{fig:training_curves}
\end{figure}

\subsection{Hyperparameter Optimisation}
\label{sec:hparam}

Hyperparameters were selected through Bayesian optimisation using
Optuna~\citep{akiba2019} with a Tree-structured Parzen Estimator (TPE)
sampler and a MedianPruner (5 startup trials, 2 warmup steps). Two
studies were run---one for the MLP baseline and one for RecurrentPPO---each
comprising 20 trials of 75\,000 timesteps with a deterministic 10-episode
evaluation to score each trial. The best configurations and study
statistics are summarised in Table~\ref{tab:bayes}.

\begin{table}[tbp]
  \centering
  \caption{Bayesian optimisation results.}
  \label{tab:bayes}
  \begin{tabular}{lcc}
    \toprule
    \textbf{Metric} & \textbf{MLP} & \textbf{RecurrentPPO} \\
    \midrule
    Completed / pruned trials & 13 / 7 & 11 / 9 \\
    Reward range & $-62.88$--$546.37$ & $27.67$--$377.07$ \\
    Mean reward & $367.14 \pm 190.40$ & $143.91 \pm 89.87$ \\
    Best learning rate & $9.31 \times 10^{-4}$ & $7.90 \times 10^{-5}$ \\
    Best $n_\text{steps}$ & 512 & 4096 \\
    Best batch size & 256 & 64 \\
    Best entropy coef & $0.01641$ & $1.96 \times 10^{-4}$ \\
    Best $\gamma$ & $0.9986$ & $0.9835$ \\
    Best clip range & $0.2320$ & $0.2320$ \\
    \bottomrule
  \end{tabular}
\end{table}

All training was conducted on a Lenovo PC with 16\,GB RAM and GPU
acceleration.

\subsection{Inference Latency}
\label{sec:latency}

Inference latency was measured across all trained policy architectures
over 1\,000 timed forward passes. Results are shown in
Figure~\ref{fig:latency} and Table~\ref{tab:latency}. The RecurrentPPO
policy produced a mean inference latency of $1.930\,\mathrm{ms}$
(p99~$=3.714\,\mathrm{ms}$, worst-case $5.310\,\mathrm{ms}$)---well
below the 10--20\,ms control period of typical upper-limb exoskeleton
platforms~\citep{supriyono2025}.

\begin{figure}[tbp]
  \centering
  \includegraphics[width=\linewidth]{fig8_latency.pdf}
  \caption{Per-step inference latency (mean~$\pm$ std) for each policy.
    Vertical lines mark the 10--20\,ms exoskeleton real-time control
    budget; all policies run well within budget.}
  \label{fig:latency}
\end{figure}

\begin{table}[tbp]
  \centering
  \caption{Inference latency across architectures (ms).}
  \label{tab:latency}
  \small
  \begin{tabular}{lrrrrrrr}
    \toprule
    \textbf{Model} & \textbf{Mean} & \textbf{Std} & \textbf{Min}
      & \textbf{p5} & \textbf{p50} & \textbf{p99} & \textbf{Max} \\
    \midrule
    Healthy MLP              & 1.267 & 0.405 & 0.840 & 0.896 & 1.090 & 2.414 & 2.799 \\
    Brady+Deg MLP            & 1.017 & 0.302 & 0.807 & 0.827 & 0.902 & 2.412 & 2.704 \\
    Deg-only MLP             & 1.056 & 0.335 & 0.808 & 0.832 & 0.913 & 2.252 & 3.275 \\
    \textbf{Brady+Deg RecPPO}& \textbf{1.930} & \textbf{0.560}
      & \textbf{1.448} & \textbf{1.488} & \textbf{1.686}
      & \textbf{3.714} & \textbf{5.310} \\
    Noisy RecPPO             & 1.851 & 0.458 & 1.447 & 1.479 & 1.666 & 3.387 & 4.110 \\
    Deg-only RecPPO          & 1.974 & 0.555 & 1.454 & 1.495 & 1.752 & 3.960 & 4.887 \\
    \bottomrule
  \end{tabular}
\end{table}

%==================================================================
\section{Experiments and Results}
\label{sec:results}

\subsection{Evaluation Methodology}

Each policy was evaluated on the same shared set of 4\,000 episode
configurations. Episodes were stratified across four severity quartiles
(Q1 mild--Q4 severe) and four radian-travelled bins reflecting the
actual distance the elbow moved during the episode:
$0.0$--$0.5$ (39.2\%), $0.5$--$1.0$ (29.2\%), $1.0$--$1.5$ (20.6\%),
and $1.5{+}$ rad (11.0\%), with 250 episodes per pairing. For every
episode, three tracks ran on identical configurations: a healthy elbow
baseline, an impaired patient with no exoskeleton assistance, and the
impaired patient with exoskeleton assistance.

The primary metric is the Pearson correlation coefficient $r$ between
the exoskeleton-assisted trajectory and the healthy baseline trajectory
over the episode. The healthy trajectory represents what the patient
\emph{intended} to do; PD breaks the link between intention and
execution. $r{=}1.0$ means the assisted movement is indistinguishable
from healthy movement; $r{=}0$ means no relationship; negative values
indicate the assisted movement opposes the intended trajectory.

Goal achievement rate---the fraction of episodes in which the target
angle was reached---is reported as a secondary metric. Mean episodic
reward is reported for all policy comparisons. All policies were trained
to maximise the same MyoSuite pose-error reward.

\subsection{Movement Restoration Performance}

The primary policy---RecurrentPPO trained on combined PD
impairments---achieved $r{=}0.494$ (std~$=0.317$) and a goal rate of
$88.8\%$ across 4\,000 stratified episodes. Relative to the
no-exoskeleton baseline ($r{=}0.223$, goal rate~$=65.1\%$), the
exoskeleton improved trajectory match by $0.271$ and goal achievement by
$23.7$ percentage points. Mean reward improved from $97.37$ to $499.73$.

Performance degraded predictably across severity quartiles:
$r{=}0.654$ at Q1, $0.522$ at Q2, $0.417$ at Q3, and $0.384$ at Q4,
with goal rates of $99.8\%$, $95.4\%$, $83.7\%$, and $76.4\%$
respectively---consistent with the exoskeleton's normalised torque
ceiling being insufficient to fully compensate at extreme impairment levels.

Performance also varied by radian travelled. Trajectory match was lowest
for smallest movements ($r{=}0.373$, 0.0--0.5 rad bin, 39.2\% of
episodes) and highest for largest movements ($r{=}0.667$, 1.5+ rad,
11.0\%). Larger movements give the LSTM more timesteps to accumulate
temporal context, improving inference of intent and impairment severity.
The 0.0--0.5 rad bin---the most common and hardest---therefore
disproportionately shapes the mean metric. Figure~\ref{fig:heatmap}
shows the full 4$\times$4 performance heatmap.

\begin{figure}[tbp]
  \centering
  \includegraphics[width=\linewidth]{fig5_heatmap.pdf}
  \caption{Mean Pearson $r$ across target-angle bins (rows) and severity
    quartiles (columns) for our combined-trained RecurrentPPO
    (\textbf{left}) versus the degeneration-only policy that represents the
    current state of the art (\textbf{right}). Both degrade toward Q4 severe
    and toward small movements (short LSTM context window), but the
    combined-trained policy retains substantially higher correlation across
    the grid.}
  \label{fig:heatmap}
\end{figure}

\subsection{Ablation 1: RecurrentPPO vs.\ MLP}

RecurrentPPO achieved $r{=}0.494$ versus the MLP's $r{=}0.439$ on the
same 4\,000 configurations, a gap of $0.055$. Goal rates were close
($88.8\%$ vs.\ $89.0\%$), but RecurrentPPO produced higher mean reward
($499.73$ vs.\ $460.29$, gap~$=39.4$)---indicating a smoother, more
proportionate trajectory toward the goal rather than merely reaching the
endpoint. The MLP stabilised at 512 timesteps; RecurrentPPO continued
varying through the full run, reflecting higher sample complexity of
learning a useful hidden state. Figure~\ref{fig:ablation1} compares the
two architectures across severity quartiles: RecurrentPPO matches or
exceeds the MLP in both trajectory match and reward at every quartile,
with the gap widening at higher severity.

\begin{figure}[tbp]
  \centering
  \includegraphics[width=\linewidth]{fig9_ablation1.pdf}
  \caption{Ablation 1: per-severity Pearson $r$ (bars) and mean episodic
    reward (lines) for the RecurrentPPO brady+deg policy versus the MLP
    brady+deg baseline, evaluated on the full brady+deg environment.
    RecurrentPPO retains higher trajectory match and reward across all
    severity quartiles.}
  \label{fig:ablation1}
\end{figure}

\subsection{Ablation 2: Combined Impairment vs.\ Degeneration-Only Training}

We claim that training on muscular degeneration alone produces a policy
that fails under the full complexity of PD. Table~\ref{tab:ablation2}
and Figure~\ref{fig:ablation_severity} test this directly.

\begin{table}[tbp]
  \centering
  \caption{Ablation 2: RecurrentPPO brady+deg vs.\ deg-only, evaluated
    on the full brady+deg environment.}
  \label{tab:ablation2}
  \begin{tabular}{lcccc}
    \toprule
    \textbf{Quartile} & \multicolumn{2}{c}{\textbf{Pearson $r$}}
      & \multicolumn{2}{c}{\textbf{Mean Reward}} \\
    \cmidrule(r){2-3}\cmidrule(l){4-5}
    & \textbf{brady+deg} & \textbf{deg-only}
    & \textbf{brady+deg} & \textbf{deg-only} \\
    \midrule
    Q1 mild    & 0.654 & 0.570 & 643.35 & 493.78 \\
    Q2         & 0.522 & 0.399 & 538.03 & 185.56 \\
    Q3         & 0.417 & 0.276 & 435.25 &  39.86 \\
    Q4 severe  & 0.384 & 0.211 & 382.30 &   9.26 \\
    \midrule
    \textbf{Overall} & \textbf{0.494} & \textbf{0.364}
      & \textbf{499.73} & \textbf{182.12} \\
    \bottomrule
  \end{tabular}
\end{table}

\begin{figure}[tbp]
  \centering
  \includegraphics[width=\linewidth]{fig4_ablation_severity.pdf}
  \caption{Ablation 2: per-severity Pearson $r$ (bars) and mean episodic
    reward (lines) for RecurrentPPO brady+deg vs.\ deg-only. At Q4 the
    degeneration-only policy's mean reward collapses to $9.26$ (vs.\
    $382.30$ for the combined policy), highlighting how it fails at severe
    bradykinesia.}
  \label{fig:ablation_severity}
\end{figure}

The degeneration-only RecurrentPPO achieved $r{=}0.364$ (std~$=0.404$),
a gap of $0.130$ from the combined-trained policy. Goal achievement fell
from $88.8\%$ to $78.9\%$. Most critically, at Q4 severe the
degeneration-only policy collapsed to $r{=}0.211$ and a mean reward of
$9.26$, compared to $r{=}0.384$ and $382.30$ for the combined policy.
The degeneration-only policy's standard deviation
grows from $0.281$ at Q1 to $0.446$ at Q4, while the combined policy
stays bounded ($0.260$--$0.341$). This pattern is consistent with the
degeneration-only agent having never learned the temporal delay pattern
of akinesia: when bradykinesia is mild the missing model matters little,
but as intensity increases the agent's inability to time torque
corrections compounds, producing both lower mean performance and higher
variance.

For MLP architectures, the same comparison showed a negligible and
reversed gap ($r{=}0.439$ brady+deg vs.\ $r{=}0.443$ deg-only). This
null result, contrasted with the $0.130$ RecurrentPPO gap, confirms that
the benefit of combined-impairment training is conditional on the
policy's capacity to exploit it. The LSTM's hidden state is what allows
combined-impairment training to translate into combined-impairment
competence.

Figure~\ref{fig:radian_performance} shows performance broken down by
radian travelled across all policies.

\begin{figure}[tbp]
  \centering
  \includegraphics[width=\linewidth]{fig6_radian_performance.pdf}
  \caption{Trajectory-tracking Pearson $r$ by radian-travelled bin for
    four policies with $\pm1\sigma$ bands. The LSTM accumulates temporal
    context over longer movements, improving performance at larger angles.
    Dashed line marks the no-exoskeleton overall baseline ($r{=}0.223$).}
  \label{fig:radian_performance}
\end{figure}

\subsection{Ablation 3: Noise Robustness}

RecurrentPPO was compared against a noise-trained variant across four
noise levels ($\sigma = 0.00, 0.01, 0.05, 0.10$) injected into muscle
activation observations only. The noise-trained policy was initially
trained for 1\,500\,000 timesteps in a clean environment, then
fine-tuned for a further 1\,500\,000 timesteps in a noisy environment
with $\sigma \sim \mathrm{Uniform}(0.01, 0.10)$ per episode.

Both policies maintained strong trajectory match across the full noise
range (Figure~\ref{fig:noise}). The clean RecurrentPPO declined by only
$\Delta r{=}0.011$ from $\sigma{=}0.00$ to $\sigma{=}0.10$. The
noise-trained variant maintained slightly higher performance, though the
gap is small. These results demonstrate that the LSTM's temporal
integration provides inherent noise robustness---the recurrent hidden
state effectively filters transient signal corruption. The tested range
reflects realistic sEMG sensing conditions: literature establishes
adequate sEMG SNR at $\geq 18\,\mathrm{dB}$~\citep{pmc8615488}.

\begin{figure}[tbp]
  \centering
  \includegraphics[width=0.85\linewidth]{fig7_noise_robustness.pdf}
  \caption{Noise robustness: clean vs.\ noisy-trained RecurrentPPO
    across observation noise $\sigma$. Performance degrades by only
    $\Delta r{=}0.011$.}
  \label{fig:noise}
\end{figure}

\begin{table}[tbp]
  \centering
  \caption{Noise robustness---Pearson $r$ and goal rate.}
  \label{tab:noise}
  \begin{tabular}{lcccc}
    \toprule
    & \multicolumn{2}{c}{\textbf{Clean RecPPO}}
    & \multicolumn{2}{c}{\textbf{Noisy-trained RecPPO}} \\
    \cmidrule(r){2-3}\cmidrule(l){4-5}
    $\boldsymbol{\sigma}$ & $r$ & Goal & $r$ & Goal \\
    \midrule
    0.00 & 0.632 & 0.906 & 0.638 & 0.923 \\
    0.01 & 0.628 & 0.897 & 0.636 & 0.895 \\
    0.05 & 0.626 & 0.903 & 0.632 & 0.897 \\
    0.10 & 0.621 & 0.899 & 0.630 & 0.897 \\
    \bottomrule
  \end{tabular}
\end{table}

\subsection{Summary}

Table~\ref{tab:summary} presents the full evaluation summary across all
policies.

\begin{table}[tbp]
  \centering
  \caption{Full evaluation summary (4\,000 shared episodes, brady+deg
    environment).}
  \label{tab:summary}
  \small
  \begin{tabular}{llccc}
    \toprule
    \textbf{Policy} & \textbf{Architecture} & \textbf{Pearson $r$}
      & \textbf{Goal Rate} & \textbf{Mean Reward} \\
    \midrule
    No exoskeleton   & ---          & 0.223 & 0.651 &  97.37 \\
    policy\_deg      & MLP          & 0.443 & 0.805 & 319.07 \\
    policy\_brady\_deg & MLP        & 0.439 & 0.890 & 460.29 \\
    policy\_deg\_recurrent & RecPPO & 0.364 & 0.789 & 182.12 \\
    \textbf{policy\_brady\_deg\_recurrent}
                     & \textbf{RecPPO} & \textbf{0.494}
                     & \textbf{0.888} & \textbf{499.73} \\
    policy\_noisy\_recurrent & RecPPO & 0.475 & 0.892 & 488.67 \\
    \bottomrule
  \end{tabular}
\end{table}

The RecurrentPPO policy trained on combined symptoms achieves the highest
trajectory match and reward. The degeneration-only RecurrentPPO performs
worst of all exoskeleton policies on the full PD environment, confirming
that temporal reasoning amplifies the consequences of training environment
mismatch.

%==================================================================
\section{Discussion}
\label{sec:discussion}

\paragraph{Interpreting results across severity.}
The RecurrentPPO policy achieved $r{=}0.494$ and $88.8\%$ goal rate
across stratified episodes covering mild (Q1, $r{=}0.654$) to severe
(Q4, $r{=}0.384$) presentations. The monotonic degradation with severity
confirms the policy is being tested across the full clinical range. For a
progressive disease, cross-severity generalisation without recalibration
is more clinically meaningful than peak accuracy on a single calibrated
patient. This matters because accessibility and adaptability---not sensor
precision---are the documented barriers to adoption: over $90\%$ of PD
patients express interest in wearable technology, yet only $8\%$ have
used a medical device~\citep{hirczy2024, reicherzer2024}.

\paragraph{Environment design as the mechanism of generalisability.}
The central result is what the ablation reveals: the degeneration-only
and combined-trained RecurrentPPO are architecturally identical, differing
only in their training environment. At Q4 severe, the degeneration-only
policy achieves $r{=}0.211$ and reward $9.26$, while the combined policy
achieves $r{=}0.384$ and reward $382.30$. This
isolates the training environment as the critical variable. The same
three-phase framework is disease-agnostic: replacing the bradykinesia and
degeneration equations with pathological models for ALS, muscular
dystrophy, or post-stroke spasticity would apply it to an entirely
different clinical population.

\subsection{Limitations}

This evaluation is entirely simulation-based---no real PD patients were
involved and no physical exoskeleton was constructed. Sim-to-real
transfer for musculoskeletal controllers remains an open
challenge~\citep{luo2025}; potential gaps in muscle dynamics, sensor
noise profiles, and body variability cannot be fully characterised without
clinical validation. However, MyoSuite was developed with validated
Hill-model dynamics and in collaboration with musculoskeletal researchers,
reducing but not eliminating this gap~\citep{caggiano2022}. The system
addresses only the elbow joint, while PD can affect gait, posture,
handwriting, and full-body coordination; the elbow remains the most
underrepresented joint in PD-specific research~\citep{capato2023}.
Finally, while the torque ceiling proved sufficient across the tested
severity range, late-stage patients at Q4 may benefit from additional
torque capacity or extended training.

\subsection{Future Work}

The immediate next step is physical exoskeleton construction and clinical
validation with real PD patients, using the cloud pipeline confirmed
viable at $1.930\,\mathrm{ms}$ mean inference latency. The three-phase
framework can be extended to hip, knee, and ankle joints for PD gait
impairment, or adapted to other neurological conditions. RLHF-based
fine-tuning using patient preference feedback on the timing and feel of
assistance would allow post-deployment personalisation without
recalibration~\citep{endrei2025}. Finally, the monotonic relationship
between severity and the policy's reward output---from $643.35$ at Q1 to
$382.30$ at Q4---suggests that exoskeleton torque demand could serve as a
continuous proxy for PD severity, enabling passive longitudinal monitoring
of disease progression through daily device use without formal UPDRS
examination.

%==================================================================
\section{Conclusion}
\label{sec:conclusion}

This paper asked whether deep reinforcement learning can produce an
autonomous elbow exoskeleton controller that provides real-time, adaptive
mitigation of Parkinson's motor impairments across diverse patients.
Three core outcomes characterise our answer. The RecurrentPPO policy
trained on combined PD symptoms achieved $r{=}0.494$ and $88.8\%$ goal
achievement across 4\,000 stratified episodes spanning the full PD
severity spectrum, without recalibration between patients. Inference
latency was $1.930\,\mathrm{ms}$ mean across 1\,000 timed forward
passes---well below the control period of typical upper-limb exoskeleton
platforms---confirming viability for real-time deployment on lightweight
hardware. Most critically, the ablation study isolates the training
environment as the mechanism of generalisability: a degeneration-only
RecurrentPPO---representative of the current state-of-the-art approach
of training exoskeletons on muscle weakness alone---collapsed to
$r{=}0.211$ and reward $9.26$ at the most severe patients, compared to
$r{=}0.384$ and $382.30$ for the combined PD-trained policy at the same
severity. This collapse, produced by architecturally identical
policies, shows that what the controller is trained on matters more than
how it is built.

ReGainX demonstrates that multi-agent RL training environments tailored
to the full complexity of disease pathology can produce exoskeleton
policies that generalise across patients without per-session calibration.
The novel three-phase framework---healthy controller training, combined PD
impairment simulation, and frozen-patient exoskeleton training---is
disease-agnostic and computationally tractable on consumer hardware,
establishing a feasible and scalable blueprint for future ML-based
assistive devices in progressive, complex neurological disease.

%==================================================================
\bibliographystyle{abbrvnat}

\begin{thebibliography}{99}

\bibitem[Agostino et al.(2015)]{agostino2015}
Agostino, R.\ et al.\ (2015).
Integrated robotics platform: $41\%$ slower speed, $29\%$ shorter
displacement off-medication.
\textit{PMC}, 6815040.
\url{https://www.ncbi.nlm.nih.gov/pmc/articles/PMC6815040/}

\bibitem[Akiba et al.(2019)]{akiba2019}
Akiba, T., Sano, S., Yanase, T., Ohta, T., and Koyama, M.\ (2019).
Optuna: A next-generation hyperparameter optimization framework.
In \textit{Proceedings of the 25th ACM SIGKDD}.
\url{https://arxiv.org/abs/1907.10902}

\bibitem[Caggiano et al.(2022)]{caggiano2022}
Caggiano, V., Wang, H., Durandau, G., Sartori, M., and Kumar, V.\ (2022).
MyoSuite: A contact-rich simulation suite for musculoskeletal motor control.
In \textit{Proceedings of Machine Learning Research}.
\url{https://proceedings.mlr.press/v168/caggiano22a.html}

\bibitem[Caligiuri et al.(2018)]{caligiuri2018}
Caligiuri, M.P.\ et al.\ (2018).
Quantitative assessment of upper limb functional impairments in PD.
\textit{Gait \& Posture}, 60, 80--85.
\url{https://www.sciencedirect.com/science/article/abs/pii/S0268003318303541}

\bibitem[Capato et al.(2023)]{capato2023}
Capato, T.T.C.\ et al.\ (2023).
Clinical assessment of upper limb impairments and functional capacity in
Parkinson's disease: a systematic review.
\textit{Arquivos de Neuro-Psiquiatria}.
\url{https://www.ncbi.nlm.nih.gov/pmc/articles/PMC10689111/}

\bibitem[Chaudhuri et al.(2025)]{chaudhuri2025}
Chaudhuri, K.R.\ et al.\ (2025).
Understanding Parkinson's disease: current trends.
\textit{Frontiers in Aging Neuroscience}.
\url{https://www.frontiersin.org/journals/aging-neuroscience/articles/10.3389/fnagi.2025.1617106/full}

\bibitem[Cools et al.(2017)]{cools2017}
Cools, A.M.\ et al.\ (2017).
Normative values and affecting factors for the elbow range of motion.
\textit{Sports Medicine}.
\url{https://www.researchgate.net/publication/319649172}

\bibitem[Dalla Gasperina et al.(2021)]{gasperina2021}
Dalla Gasperina, S.\ et al.\ (2021).
Review on patient-cooperative control strategies for upper-limb
rehabilitation exoskeletons.
\textit{Frontiers in Robotics and AI}.
\url{https://www.ncbi.nlm.nih.gov/pmc/articles/PMC8688994/}

\bibitem[Durandau et al.(2019)]{durandau2019}
Durandau, G.\ et al.\ (2019).
Neuromusculoskeletal controller for exoskeleton assistance.
Referenced via MyoSuite documentation~\citep{caggiano2022}.

\bibitem[Endrei et al.(2025)]{endrei2025}
Endrei, T.\ et al.\ (2025).
Learning to suppress tremors: a DRL-enabled soft exoskeleton for
Parkinson's patients.
\textit{Frontiers in Robotics and AI}.
\url{https://www.frontiersin.org/journals/robotics-and-ai/articles/10.3389/frobt.2025.1537470}

\bibitem[Fortunati et al.(2024)]{fortunati2024}
Fortunati, M.\ et al.\ (2024).
Lower-limb exoskeletons for gait training in Parkinson's disease:
state of the art.
\textit{Healthcare}.
\url{https://pmc.ncbi.nlm.nih.gov/articles/PMC11353983/}

\bibitem[Hart et al.(2023)]{hart2023}
Hart, N.\ et al.\ (2023).
Prevalence of sarcopenia in Parkinson's disease: systematic review.
\textit{Neurological Sciences}.
\url{https://www.ncbi.nlm.nih.gov/pmc/articles/PMC10641055/}

\bibitem[Hirczy et al.(2024)]{hirczy2024}
Hirczy, S.\ et al.\ (2024).
The current state of wearable device use in Parkinson's disease: a survey.
\url{https://www.ncbi.nlm.nih.gov/pmc/articles/PMC11701158/}

\bibitem[Ingram et al.(2021)]{ingram2021}
Ingram, L.A.\ et al.\ (2021).
Quantifying upper limb motor impairment in people with Parkinson's disease.
\textit{PeerJ}.
\url{https://pubmed.ncbi.nlm.nih.gov/33604177/}

\bibitem[Luo et al.(2023)]{luo2023}
Luo, S.\ et al.\ (2023).
Robust walking control of a lower limb rehabilitation exoskeleton via DRL.
\textit{Journal of NeuroEngineering and Rehabilitation}.
\url{https://link.springer.com/article/10.1186/s12984-023-01147-2}

\bibitem[Luo et al.(2024)]{luo2024}
Luo, S.\ et al.\ (2024).
Experiment-free exoskeleton assistance via learning in simulation.
\textit{Nature}.
\url{https://www.nature.com/articles/s41586-024-07382-4}

\bibitem[Luo et al.(2025)]{luo2025}
Luo, S.\ et al.\ (2025).
Robust neural network controller for a lower limb exoskeleton with
minimal sensors.
\textit{Journal of NeuroEngineering and Rehabilitation}.
\url{https://link.springer.com/article/10.1186/s12984-025-01854-y}

\bibitem[MDPI Sensors(2025)]{mdpi2025}
Supriyono, Dragusanu, and Malvezzi\ (2025).
A comprehensive review of elbow exoskeletons.
\textit{Sensors}.
\url{https://www.mdpi.com/1424-8220/25/14/4263}

\bibitem[Morrey et al.(1981)]{morrey1981}
Morrey, B.F.\ et al.\ (1981).
Functional range of motion of the elbow.
\textit{Journal of Bone and Joint Surgery}.
\url{https://pubmed.ncbi.nlm.nih.gov/7775772/}

\bibitem[Myoelectric(2022)]{myoelectric2022}
PMC\ (2022).
Myoelectric control systems for upper limb wearable robotic exoskeletons.
\textit{Sensors}.
\url{https://pmc.ncbi.nlm.nih.gov/articles/PMC9655258/}

\bibitem[Parkinson's Foundation(2026)]{parkinson2026}
Parkinson's Foundation\ (2026).
\$82.2 billion economic burden.
\url{https://www.parkinson.org/blog/research/economic-burden-2024}

\bibitem[PMC(2024)]{reicherzer2024}
Reicherzer, L.\ et al.\ (2024).
Barriers and facilitators to the use of wearable robots as assistive
devices.
\url{https://pmc.ncbi.nlm.nih.gov/articles/PMC11327840/}

\bibitem[PMC(2025a)]{pmc11852623}
PMC\ (2025).
LSTM-enabled EMG-controlled upper arm exoskeleton.
\url{https://pmc.ncbi.nlm.nih.gov/articles/PMC11852623/}

\bibitem[PMC(2025b)]{pmc12451468}
Carnevale, A.\ et al.\ (2025).
Portable exoskeletons for upper limb rehabilitation: a systematic review.
\url{https://pmc.ncbi.nlm.nih.gov/articles/PMC12451468/}

\bibitem[PMC(2025c)]{pmc8615488}
PMC\ (2021).
An ultra-low power surface EMG sensor for wearable applications.
\url{https://pmc.ncbi.nlm.nih.gov/articles/PMC8615488/}

\bibitem[Ponsoni et al.(2023)]{ponsoni2023}
Ponsoni, A.\ et al.\ (2023).
Prevalence of sarcopenia in PD: meta-analysis.
\url{https://pubmed.ncbi.nlm.nih.gov/36413812/}

\bibitem[Raethjen et al.(2018)]{raethjen2018}
Raethjen, J.\ et al.\ (2018).
Quantitative analysis of bradykinesia and rigidity in PD.
\url{https://pmc.ncbi.nlm.nih.gov/articles/PMC5853013/}

\bibitem[Raffin et al.(2021)]{sb3}
Raffin, A.\ et al.\ (2021).
Stable-Baselines3: Reliable reinforcement learning implementations.
\textit{JMLR}.
\url{https://jmlr.org/papers/v22/20-1364.html}

\bibitem[Raffin and Kanervisto(2023)]{sb3contrib}
Raffin, A.\ and Kanervisto, A.\ (2023).
SB3-Contrib: Experimental RL code for Stable Baselines3.
\url{https://github.com/Stable-Baselines3/stable-baselines3-contrib}

\bibitem[Ramesh and Bilal(2022)]{ramesh2022}
Ramesh, V.\ and Bilal, E.\ (2022).
Detecting motor symptom fluctuations in Parkinson's disease with
generative adversarial networks.
\textit{npj Digital Medicine}.
\url{https://www.nature.com/articles/s41746-022-00674-x}

\bibitem[ResearchGate(2023)]{researchgate2023}
ResearchGate discussion\ (2023).
How can I find a dataset of pathological EMG signals?
\url{https://www.researchgate.net/post/How_can_I_find_a_dataset_of_pathological_EMG_signals}

\bibitem[arXiv(2025)]{arxiv2025}
arXiv\ (2025).
AI-enabled detection and assessment of PD using multimodal data: a survey.
\url{https://arxiv.org/html/2502.10703v1}

\bibitem[Siciliano et al.(2024)]{siciliano2024}
Siciliano, M.\ et al.\ (2024).
Fatigue in Parkinson's disease: pooled prevalence $\approx$50\%.
\url{https://pmc.ncbi.nlm.nih.gov/articles/PMC11626731/}

\bibitem[Supriyono et al.(2025)]{supriyono2025}
Supriyono, Dragusanu, M., and Malvezzi, M.\ (2025).
A comprehensive review of elbow exoskeletons: classification by structure,
actuation, and sensing technologies.
\textit{Sensors}.
\url{https://pmc.ncbi.nlm.nih.gov/articles/PMC12299220/}

\bibitem[Wang et al.(2022)]{wang2022}
Wang, H., Caggiano, V.\ et al.\ (2022).
MyoSim: Fast and physiologically realistic MuJoCo models for
musculoskeletal and exoskeletal studies.
In \textit{IEEE ICRA}.
\url{https://dl.acm.org/doi/10.1109/ICRA46639.2022.9811684}

\bibitem[Zajac(1989)]{zajac1989}
Zajac, F.E.\ (1989).
Muscle and tendon: properties, models, scaling, and application to
biomechanics and motor control.
\textit{Critical Reviews in Biomedical Engineering}, 17(4), 359--411.
\url{https://pubmed.ncbi.nlm.nih.gov/2676342/}

\end{thebibliography}

\end{document}
